% stub

std::variant 和 std::optional 被称为“词汇（vocabulary）类型”，因为你可以通过它们传达更多的设计信息。你的代码可以变得更加紧凑、更具表达力。本章将通过一个示例展示 std::optional 和 std::variant 如何帮助你重构一个函数。我们会从一些遗留代码开始，经过若干步骤，最终得到一个好得多的方案。为了帮助你更好地理解，我们将看到每一步的优缺点。

\section{使用场景}

设想一个函数，它接收某款游戏中当前的鼠标选区。该函数扫描所选范围，并计算出若干输出：

\begin{itemize}
\item 正在执行动画的对象数量
\item 选区中是否包含任何民用单位
\item 选区中是否包含任何战斗单位
\end{itemize}

现有代码看起来像这样：

\begin{codebox}[cpp]
class ObjSelection
{
public:
    bool IsValid() const { return true; }
    // 更多代码...
};

bool CheckSelectionVer1(const ObjSelection &objList,
                        bool *pOutAnyCivilUnits,
                        bool *pOutAnyCombatUnits,
                        int *pOutNumAnimating);
\end{codebox}

如你所见，该函数使用了大量输出参数（以裸指针的形式），并返回 true/false 来指示成功与否（例如，输入的选区可能无效）。

该函数具体的实现现在并不重要，但下面是一个调用该函数的示例代码：

\begin{codebox}[cpp]
ObjSelection sel;
bool anyCivilUnits { false };
bool anyCombatUnits { false };
int numAnimating { 0 };
if (CheckSelectionVer1(sel, &anyCivilUnits, &anyCombatUnits, &numAnimating))
{
    // ...
}
\end{codebox}

我们该如何改进这个函数呢？可能要从以下几个方面入手：

\begin{itemize}
\item 看看调用方的代码：我们必须创建所有用来保存输出的变量。如果你在多处调用这个函数，无疑会产生代码重复。
\item 输出参数：Core Guidelines（核心指南）建议不要使用它们。
  \begin{itemize}
  \item F.20：对于“out”输出值，优先使用返回值而非输出参数\footnote{\url{https://isocpp.github.io/CppCoreGuidelines/CppCoreGuidelines\#f20-for-out-output-values-prefer-return-values-to-output-parameters}}。
  \end{itemize}
\item 如果使用裸指针，你就必须检查它们是否有效。如果改用引用作为输出参数，则可以省去这些检查。
\item 那扩展这个函数呢？如果你需要再增加一个输出参数怎么办？
\end{itemize}

还有别的吗？

你会如何重构它呢？

受 Core Guidelines 与 C++17 新特性的启发，下面是我们改进这段代码的方法：

\begin{enumerate}
\item 把输出参数重构为一个将被返回的 tuple。
\item 把 tuple 重构为一个独立的结构体，并将 tuple 缩减为一个 pair。
\item 使用 std::optional 来表达该值是否被计算出来。
\item 使用 std::variant 不仅传达可选的结果，还传达完整的错误信息。
\end{enumerate}

\section{Tuple 版本}

第一步是把输出参数转换成一个 tuple，并从函数返回它。

根据 F.21：要返回多个“out”值，最好返回一个 tuple 或 struct\footnote{\url{https://isocpp.github.io/CppCoreGuidelines/CppCoreGuidelines\#f21-to-return-multiple-out-values-prefer-returning-a-tuple-or-struct}}：

\begin{myQuote}
返回值是一种“仅输出”的值，本身具有自说明性。请注意，C++ 确实支持多返回值，其约定方式是使用 tuple（包括 pair），在调用处还可以借助 tie 带来额外的便利。
\end{myQuote}

改动之后，代码可能变成这样：

\begin{codebox}[cpp]
std::tuple<bool, bool, bool, int>
CheckSelectionVer2(const ObjSelection &objList)
{
    if (!objList.IsValid())
        return {false, false, false, 0};
    // 局部变量：
    int numCivilUnits = 0;
    int numCombat = 0;
    int numAnimating = 0;
    // 扫描...
    return {true, numCivilUnits > 0, numCombat > 0, numAnimating };
}
\end{codebox}

好一点了……不是吗？

\begin{itemize}
\item 无需检查裸指针
\item 代码相当有表达力
\end{itemize}

此外，在调用处，你可以使用结构化绑定（Structured Bindings）来接收返回的 tuple：

\begin{codebox}[cpp]
auto [ok, anyCivil, anyCombat, numAnim] = CheckSelectionVer2(sel);
if (ok)
{
    // ...
}
\end{codebox}

遗憾的是，这个版本未必是最好的。例如，存在忘记 tuple 中各输出顺序的风险。

函数扩展的问题也依然存在。因此，当你想增加另一个输出值时，就必须扩展这个 tuple 以及调用处。

我们可以通过再前进一步来解决这个问题：不再使用 tuple，而是使用一个结构体（Core Guidelines 也建议这么做）。

\section{一个独立的结构体}

这些输出看起来表示的是相关的数据。因此，把它们封装到一个名为 SelectionData 的 struct 中或许是个好主意。

\begin{codebox}[cpp]
struct SelectionData
{
    bool anyCivilUnits { false };
    bool anyCombatUnits { false };
    int numAnimating { 0 };
};
\end{codebox}

然后你可以把函数改写为：

\begin{codebox}[cpp]
std::pair<bool, SelectionData> CheckSelectionVer3(const ObjSelection &objList)
{
    SelectionData out;
    if (!objList.IsValid())
        return {false, out};
    // 扫描...
    return {true, out};
}
\end{codebox}

而调用处：

\begin{codebox}[cpp]
if (auto [ok, selData] = CheckSelectionVer3(sel); ok)
{
    // ...
}
\end{codebox}

代码使用了 std::pair，所以我们仍然保留了成功标志，它不属于那个新结构体。

我们在这里获得的主要优势是更好的代码结构与可扩展性。如果你想添加一个新参数，只需扩展该结构体。而在以前——函数的声明里是一长串输出参数——你就得改动多得多的代码。

但 std::pair<bool, MyType> 难道不类似于 std::optional 的概念吗？

\section{使用 std::optional}

引自 std::optional 那一章：

std::optional 是一种包装类型，用来表达“可空（null-able）”类型。它要么包含一个值，要么为空。它不会使用任何额外的内存分配。

这似乎非常适合我们的代码。我们可以去掉 ok 变量，转而依赖 optional 的语义。

新版本的代码：

\begin{codebox}[cpp]
std::optional<SelectionData> CheckSelectionVer4(const ObjSelection &objList)
{
    if (!objList.IsValid())
        return std::nullopt;
    SelectionData out;
    // 扫描...
    return out;
}
\end{codebox}

而调用处：

\begin{codebox}[cpp]
if (auto ret = CheckSelectionVer4(sel); ret.has_value())
{
    // 通过 *ret 访问，或者甚至用 ret->
    // ret->numAnimating
}
\end{codebox}

optional 版本有哪些优势呢？

\begin{itemize}
\item 形式简洁且富有表达力——optional 表达的是可空类型。
\item 高效——optional 的实现不允许使用额外的存储（例如动态内存）。所包含的值应当分配在 optional 存储中、针对类型 T 做了适当对齐的区域里。
\end{itemize}

\section{使用 std::variant}

前面使用 std::optional 的最后一个实现遗漏了一个关键方面：错误处理。我们无从得知某个值没有被计算的原因。例如在使用 std::pair 的版本中，我们可以返回一个错误码来表明原因。对此我们能做什么呢？

如果你需要关于函数内部可能出错的全部信息，可以考虑改用 std::variant 的另一种方案。

\begin{codebox}[cpp]
enum class [[nodiscard]] ErrorCode
{
    InvalidSelection,
    Undefined
};

variant<SelectionData, ErrorCode> CheckSelectionVer5(const ObjSelection &objList)
{
    if (!objList.IsValid())
        return ErrorCode::InvalidSelection;
    SelectionData out;
    // 扫描...
    return out;
}
\end{codebox}

如你所见，代码使用了 std::variant，它有两种可能的备选项：要么是 SelectionData，要么是 ErrorCode。它几乎就像一个 pair，只不过你始终只会看到其中一个活跃的值。

你可以这样使用上面的实现：

\begin{codebox}[cpp]
if (auto retV5 = CheckSelectionVer5(sel);
    std::holds_alternative<SelectionData>(retV5))
{
    std::cout << "ok..."
              << std::get<SelectionData>(retV5).numAnimating << '\n';
}
else
{
    switch (std::get<ErrorCode>(retV5))
    {
    case ErrorCode::InvalidSelection:
        std::cerr << "Invalid Selection!\n";
        break;
    case ErrorCode::Undefined:
        std::cerr << "Undefined Error!\n";
        break;
    }
}
\end{codebox}

正如你所看到的，使用 std::variant 能比使用 std::optional 获得更多信息。你可以返回错误码，并对可能出现的失败做出响应。

std::variant<ValueType, ErrorCode> 可能是 std::expected 的一种实现——std::expected 是一个新的词汇类型，未来或许会被纳入标准库。

\section{小结}

你可以在以下位置试验这段代码：\texttt{Chapter Refactoring With Optional And Variant/refactoring\_optional\_variant.cpp}。

在本章中，你已经看到如何把一堆难看的输出参数重构为更优雅的 std::optional 版本。这个 optional 包装清晰地表达了计算所得的值可能不存在。此外，你还看到了如何把多个函数参数封装到一个独立的结构体中。拥有这样一个独立的类型，可以让你轻松扩展代码，同时保持逻辑结构不变。

最后，如果你需要函数内部错误的完整信息，那么也可以考虑改用 std::variant。这种类型让你可以返回一个完整的错误码。
